\documentclass[10pt,a4paper]{amsart}
\usepackage[margin=19mm]{geometry}
\usepackage{amsmath,amssymb,mathtools}
\usepackage[scr=rsfs,cal=euler]{mathalfa}
\usepackage{xfrac}
\usepackage[unicode,hidelinks,bookmarks=false]{hyperref}
\newcommand{\ie}{i.e.\ }
\newcommand{\cf}{cf.\ }
\newcommand{\resp}{respectively\ }
\newcommand{\Subsection}{\subsection}
\providecommand{\nobreakdash}{\nobreak\hbox{-}}
\setlength{\emergencystretch}{2em}
\multlinegap0pt
\usepackage{bm}
\usepackage{bbm}
\usepackage{dsfont}
\usepackage{xspace}
\usepackage{cancel}
\usepackage{mathtools}
\usepackage{amsthm}
\makeatletter
\@ifundefined{theorem}{\newtheorem{theorem}{Theorem}}{}
\@ifundefined{proposition}{\newtheorem{proposition}[theorem]{Proposition}}{}
\makeatother
\DeclareMathOperator{\h}{\mathbf{h}} % Metric functional
\title[Subinvariant metric functionals for nonexpansive mappings]{Subinvariant metric functionals for nonexpansive mappings}
\author{Armando W. Guti\'errez, Olavi Nevanlinna}
\date{Source: arXiv:2407.04234v1 (CC BY 4.0)}
\begin{document}
\maketitle

\noindent\emph{Source and licence.} Subinvariant metric functionals for nonexpansive mappings. Version arXiv:2407.04234v1 (CC BY 4.0), distributed under the Creative Commons Attribution 4.0 International licence, as recorded in the arXiv licence metadata for this exact version and shown on the abstract page https://arxiv.org/abs/2407.04234v1. Full rights notice: licenses/arxiv-2407.04234-CC-BY-4.0.txt.\par\smallskip

\noindent\textit{Editorial scope.} Counted unit: the Proposition whose source label is \texttt{None} in the article (source0.tex lines 635--653, source label \texttt{None}), delivered together with the article's own complete original proof of it (source0.tex lines 655--705, one \texttt{proof} environment of 51 lines). No mathematics is written, repaired or re-derived: statement and proof are the article's own text line for line. Required context retained uncounted: none. Every definition, display and label of those lines is retained unchanged. Omitted from this package and not redistributed: 1--634, 654--654, 706--1383. Nothing inside a retained excerpt is omitted or altered: every retained body is delivered exactly as the article's lines are written, so no omission receipt of any kind accompanies this package. Citations: the retained excerpts carry no citation command at all, so no bibliography entry of the article is transcribed and no reference is redistributed. Reference adaptation: none was needed and none was made; every reference inside the delivered unit resolves inside it, so no unresolved reference or citation is produced. Printed slips kept literal: no printed word, symbol, hypothesis or number of the counted statement or of its proof was altered. Layout and class: the article's own class is \texttt{article} with packages including amsmath, amsthm, amssymb, mdframed; the wrapper is the lane's pinned \texttt{amsart} editorial wrapper at 10pt on A4, which restates the theorem-like environments the retained text needs and adds the article's own macro definitions that the retained text needs (\texttt{\textbackslash h}), with extra wrapper packages (bm, bbm, dsfont, xspace, cancel, mathtools). The delivered numbering is the unit's own. Assets: no retained span contains a figure environment, a graphics-inclusion command, a PDF-page inclusion, a table or a TikZ picture; no image file is copied into this package, the package declares no asset dependency and no image file is redistributed with it. Licence: version arXiv:2407.04234v1 of this article is distributed under the Creative Commons Attribution 4.0 International licence, as recorded in the arXiv licence metadata for this exact version and shown on the abstract page https://arxiv.org/abs/2407.04234v1. A full rights notice is delivered at licenses/arxiv-2407.04234-CC-BY-4.0.txt. Characters: every retained excerpt is pure ASCII, so the delivered engine, XeLaTeX with its default Latin Modern text font, reports no missing character.
\begin{proposition}  
For the metric space $X_1$ the following 
properties hold:
\begin{enumerate}
\item We have 
$\,\h \in (X_1)^\diamondsuit\,$ if and only 
if there are two metric functionals 
$\,\h^A\in A^\diamondsuit\,$ and 
$\,\h^B\in B^\diamondsuit$ such that
$$
\h(x) = \h^A(a) + \h^B(b)\ \ 
\text{for all}\ \ x=(a,b)\in X_1. 
$$
\item The metric space $X_1$ has ZFP if and 
only if at least one of $A$ and $B$ has ZFP.
\item The metric space $X_1$ has UMP if and 
only if both $A$ and $B$ have UMP.
\end{enumerate}
\end{proposition}

\begin{proof}
Let us fix two points $a_0\in A$ and 
$b_0\in B$. 

The first statement follows from the 
formula
$$d_1((a,b),(w,z)) - d_1((a_0,b_0),(w,z)) 
    = \h_{w}^{A}(a)+ \h_{z}^{B}(b),$$ 
where $\h_{w}^{A}(a)=d_A(a,w)-d_A(a_0,w)$ 
and $\h_{z}^{B}(b)=d_B(b,z)-d_B(b_0,z)$.     

Let us consider the second statement. 
If both $A$ and $B$ have identically 
vanishing metric functionals, their sum 
also vanishes. Now, let us assume that $A$ 
has ZFP and consider a metric functional 
$\h = \h^A + \h^B$. We know that there exists
$\hat{a} \in A$ such that 
$\h^A(\hat{a})\neq 0$. Since $\h^B(b_0)=0$, 
we have 
$\h(\hat{a},b_0) = \h^A(\hat{a}) \neq 0$.
If $B$ has ZFP instead, then we have
$\h(a_0,\hat{b})\neq 0$ for some 
$\hat{b}\in B$.    

To show the last statement, let us assume 
first that both $A$ and $B$ have UMP. Then, 
for every metric functional $\h = \h^A + \h^B$
there are $\hat{a}\in A$ and $\hat{b}\in B$
such that $\h^A(\hat{a}) < \h^A(a)$ for all 
$a\neq \hat{a}$ and $\h^B(\hat{b}) < \h^B(b)$ 
for all $b\neq \hat{b}$. This implies that
$\h(\hat{a},\hat{b}) < \h(a,b)$ for all 
$(a,b) \neq (\hat{a},\hat{b})$. Now, let us 
assume that $X_1$ has UMP and consider two 
metric functionals $\h^A\in A^\diamondsuit$ 
and $\h^B\in B^\diamondsuit$. Since $X_1$ has
UMP, there are two points $\hat{a}\in A$ and 
$\hat{b}\in B$ such that 
$$
\h^A(\hat{a})+\h^B(\hat{b}) 
< \h^A(a)+\h^B(b)\ \ \text{for all}\ \ 
(a, b)\neq (\hat{a}, \hat{b}). 
$$
In particular, by evaluating the inequality
shown above at the point $(a,\hat{b})$ with 
$a\neq \hat{a}$, we have 
$\h^A(\hat{a}) < \h^A(a)$. If we consider
the point $(\hat{a},b)$ with $b\neq \hat{b}$ 
instead, we have $\h^B(\hat{b}) < \h^B(b)$.
\end{proof}

\end{document}
